\section{Item (1): the regularisation hypothesis \texorpdfstring{$\Hreg$}{(H\_reg)} for the bubbles}\label{sec:hreg}
\paragraph{Inner degrees.} Consider a bubble graph with cubic vertices at $x$ and $x'$. Let $x\to W$ with $x'$ fixed away from $W$. Near $W$ each propagator leg at $x$ (two internal, one external) is a sum of jets $r^{m}e^{in\varphi}$, and at most $k\le2$ transverse derivatives act at the vertex. The integrand is therefore a finite sum of terms $r^{s}e^{iN\varphi}\times$(smooth in the longitudinal variables) with
\[ s=m_1+m_2+m_3-k,\qquad N=n_1+n_2+n_3, \]
and only $N=0$ survives the angular integral. The term is locally integrable in the 2-dimensional transverse plane iff $s>-2$. Terms built only from $n=0$ jets are products of derivatives of smooth functions ($r^{|\sigma|}e^{i\sigma\varphi}$ in the helicity frame are polynomials), hence smooth; they cause no apex singularity.
\begin{theorem}[Inner-degree classification]\label{thm:inner}
Computed exhaustively in \texttt{code/inner\_degrees.py} and \texttt{code/inner\_degrees\_q24.py} for the vertices of D01 $(h,h|h)$, D03 $(c,c|h)$, D04 $(h,\phi|\phi)$ and $(\phi,h|\phi)$, D11 $(\phi,\phi|h)$ and D13 $(c,h|K)$, $(c,h|c)$ (D06 is the D01$+$D03 channel), with $|n|\le5$, $j\le1$, on the grid $\{k/97\}\cup\{p/q:\ q\le24\}$ (276 values), both branches:
\begin{enumerate}
\item[(i)] On CR, a non-integrable inner degree ($s\le-2$, $N=0$, some $n_i\ne0$) that is an \emph{integer} occurs exactly at $\alpha\in\{1/2,\,2/3,\,4/5\}$. The affected graphs are D01 (6 terms at $1/2$, 42 at $2/3$), D03 (4 at $2/3$), D13 ($2+0$ at $1/2$, $16+4$ at $2/3$), and at $4/5$ D01, D03 and D13 (example: $n=(1,1,-2)$, $s=4t-6-k=-2$ at $k=1$).
\item[(ii)] On F this happens only at $\alpha=1/2$ (where F and CR coincide).
\item[(iii)] D04 and D11 have no non-integrable inner degree with $n\ne0$ at any $\alpha$; they are \emph{inner-clean}.
\item[(iv)] At every other $\alpha$ all non-integrable inner degrees are non-integer. At $\alpha=7071/10000$ the minima are $s_{\min}=-3.1715$ (D01, D13a) and $-2.1715$ (D03, D13b).
\end{enumerate}
The grid is complete. For $\alpha<1/2$ there are no flips, all $m\ge0$, and an integer $s\le-2$ needs $k=2$ and three $m=0$ jets, one with $n\ne0$, i.e.\ $\nu=|nt+\sigma|=0$, impossible for $t>2$. At $\alpha=1/2$ ($t=2$), $\nu=0$ at $(\sigma,n)=(\pm2,\mp1)$ gives the non-smooth $r^0e^{\mp i\varphi}$. For $\alpha\in(1/2,1)$, each flip has $n=\mp1$ and $m\ge t-2>-1$; three flips cannot satisfy $N=0$, so at most two occur and contribute $>-2$. Since $k\le2$, $s\le-2$ forces every non-flipped jet to have $m<2$, and $m\ge|n|t-2$ gives $|n|<4/t<4$. Hence $s=At+B$ with $|A|\le9$, and $s\in\mathbb Z$ forces $\alpha=|A|/K$ with $|A|<K<2|A|\le18$. All such values lie in the grid.
\end{theorem}
The degree count is power counting: vertex tensor contractions were not evaluated. They can only remove terms, so the lists in (i) are upper bounds and (iii), (iv) are unaffected.
\begin{proposition}[$\Hreg$ at inner-non-resonant $\alpha$, CR; modulo (J)]\label{prop:hreg}
Let $\alpha\notin\{1/2,2/3,4/5\}$ on CR. Define the bubble regularisation as: longitudinal dimensional continuation (XVI-P DRED, $d-2$ longitudinal dimensions, transverse plane exactly 2-dimensional), composed with, at each vertex approaching $W$ alone, the \emph{unique homogeneous extension} of every non-integrable inner term $r^se^{iN\varphi}$ (unique because $s\notin\{-2,-3,\dots\}$; Hörmander, Thm.~3.2.3). Assume
\begin{itemize}
\item[(J)] the joint region $x,x'\to W$ is regulated by $d$, i.e.\ its joint-scaling Mellin exponent depends on $d$ non-degenerately.
\end{itemize}
Then the regularised one-loop $\ka^2$ two-point functional satisfies the linearised Slavnov--Taylor identity up to terms $(d-4)\times$(regularised), so its pole part is $\sigma_0$-closed. Hence (v2 Thm.~6.1) the 26 Ward-constrained pole coefficients vanish (27 at $1/3$).
\end{proposition}
\begin{proof}[Proof sketch]
(a) \emph{Domain.} On CR the BRST variations map lattice jets into lattice jets ($K_2c$ has no foreign atom in any sample; v1/v2).
(b) \emph{Commutation.} The unique homogeneous extension $E_s$ of degree $s$ commutes with $\partial_i$ ($E_{s-1}\partial=\partial E_s$, both are homogeneous extensions of $\partial f$ and $s-1$ is again non-exceptional) and with multiplication by $x_i$ (degree $s+1$, non-exceptional). The Ward identities of the propagators hold as distributions including $W$ (v1). Hence every algebraic step of the flat ST derivation goes through for the extended integrands.
(c) \emph{Longitudinal DR} commutes with longitudinal derivatives and translations. The only breaking is the mismatch between quasi-$d$ momenta and quasi-4 indices, which is evanescent, $(d-4)\times$ an object with at most a simple pole, hence finite at one loop.
(d) By (J) the joint region produces a pole in $d$ only, and it inherits (a)--(c). Finite apex breakings are WZ-closed and exact by $H^1_{(2)}=0$, so they are absorbed by finite inventory counterterms.
\end{proof}
\paragraph{Status and obstruction.} (J) is argued, not proved. A proof needs the joint Mellin representation of \S\ref{sec:bubbles} and a check that its exponent is not $d$-independent for some channel. At $\alpha\in\{1/2,2/3,4/5\}$ the obstruction is precise. The listed integer terms with $N=0$ have non-vanishing angular moment ($\int d\varphi\ne0$ at $s=-2$), so there is no homogeneous extension: they are genuine transverse logarithms that longitudinal DR does not regulate. They need either an extra regulator (e.g.\ analytic $r^{2\epsilon}$, whose BRST behaviour is not established) or cubic apex counterterms on $W$. If the breaking induced by such a regulator is local on $W$ and WZ-consistent, it is removable when the three-field cohomology vanishes at that $\alpha$. We computed $H^1_{(3)}=0$ at $1/2$ and $2/3$ (\S\ref{sec:tri}), and at $4/5$ (all-jet run, candidate counts generic; \S\ref{sec:tri}). That gives a conditional route, not a proof. On F, $\Hreg$ fails already for the $E=1$ row (v2), so no F statement is made.
\textbf{Item (1) outcome: partially earned}: proved modulo (J) and standard lemmas at inner-non-resonant $\alpha$; obstructed (precisely located) at $\{1/2,2/3,4/5\}$.

\section{Item (2): apex poles of the bubbles}\label{sec:bubbles}
\paragraph{Reduction.} Let $B(x,x')$ be a bubble integrand after the inner extensions of \S\ref{sec:hreg}. Write $x=(\rho\omega,y)$, $x'=(\rho\omega',y')$ with $(\omega,\omega')$ on the unit scaling sphere $\Sigma$ of the transverse pair $(x_\perp,x'_\perp)\in\mathbb R^2\times\mathbb R^2$ (dimension 3). Then the apex pole of the bubble is
\[ \mathrm{Pole}_{d=4}\,\Gamma_W=\frac{1}{4-d}\,\mathrm{Res}_{z=z_*(d)}\int_\Sigma d\omega\,d\omega'\;\mathcal M[B](z;\omega,\omega';p_\parallel), \]
where $\mathcal M$ is the Mellin transform in $\rho$, taken after the bulk Hadamard subtraction at $x=x'$ (which reproduces the flat 21-direction pole of Theorem~v1) and after the inner extensions. Projected onto the 34-dimensional basis (equivalently, by v2 Thm.~6.1 and Prop.~\ref{prop:hreg}, onto the 8 closed directions), it gives the 6 exact and the 2 $\phi\phi$ values. For D04 one can do better. Split the graviton propagator as $\Pi\,G_0+\sum_S\Pi_S(G_{b(S)}-G_0)$. The first part obeys the flat DRED identity $V\Pi V'=2p^2(l-p)^2$ and reduces to a tadpole (route of \S\ref{sec:tad}). The second part has a bulk-smooth kernel $G_{b(S)}-G_0$, so its pole is purely apex, with no Hadamard subtraction. D04 and D11 need no inner extension at any $\alpha$ (Thm.~\ref{thm:inner}(iii)).
\paragraph{Obstruction.} The residue is a 3-dimensional singular integral over $\Sigma$ of products of the exact Bernoulli kernels (and of their first and second derivatives), with an extra longitudinal loop integral in $d-2$ dimensions. We did not evaluate it, for D04 alone or D06 alone or any other bubble. The Synge (coincidence) data used for tadpoles do not suffice, because $G_{b}-G_0$ is not smooth at $r=r'=0$. The I-B-VP-1 technology handles one-point extensions ($U_L$ of $r^{-2}$, ambiguity $\delta_W$) and is used in the inner step. It does not provide the joint two-point residue.
\textbf{Item (2) outcome: open}, with the reduction above. The 26 Ward-constrained combinations vanish only modulo (J) (\S\ref{sec:hreg}). The 6 exact and 2 $\phi\phi$ values are \emph{not} computed (missing $\ne$ zero).
